% !TEX program = pdflatex
% Para gerar o GABARITO, tire o % da linha abaixo e compile.
% \printanswers
\documentclass[11pt,a4paper,addpoints]{exam}
\usepackage[T1]{fontenc}
\usepackage[utf8]{inputenc}
\usepackage[brazilian]{babel}
\usepackage{lmodern}
\usepackage{amsmath,amssymb}
\usepackage[margin=2.2cm]{geometry}
\usepackage{xcolor}
\definecolor{tocaHoney}{HTML}{C47E12}

% Textos da classe exam em português.
\pointpoints{ponto}{pontos}
\bonuspointpoints{ponto extra}{pontos extras}
\renewcommand{\solutiontitle}{\noindent\textbf{Resposta:}\enspace}
\hpword{Pontos:}
\vpword{Pontos}
\hqword{Questão:}
\vqword{Questão}
\htword{Total}
\vtword{Total}
\hsword{Nota:}
\vsword{Nota}
\shadedsolutions
\definecolor{SolutionColor}{HTML}{FBF1DF}
\qformat{\textbf{Questão \thequestion} \hfill (\totalpoints\ \points)}

\pagestyle{headandfoot}
\firstpageheader{{{INSTITUTION}}}{}{{{DATE}}}
\runningheader{{{TITLE}}}{}{Página \thepage\ de \numpages}
\firstpagefooter{}{}{}
\runningfooter{}{}{}

\begin{document}

\begin{center}
  {\Large\bfseries {{TITLE}}}\par\medskip
  Professor(a): {{AUTHOR}}\par
  {\color{tocaHoney}\rule{3cm}{2pt}}
\end{center}

\vspace{0.5em}
\noindent\makebox[0.65\textwidth]{Nome:\enspace\hrulefill}\hfill\makebox[0.3\textwidth]{Matrícula:\enspace\hrulefill}

\vspace{1em}
\begin{center}
  \gradetable[h][questions]
\end{center}

\noindent\textbf{Instruções.} A prova vale \numpoints\ pontos e tem \numquestions\ questões. Justifique as respostas. Duração: 2 horas.

\begin{questions}

\question[2] Defina custo de oportunidade e dê um exemplo.
\begin{solution}
  O valor da melhor alternativa de que se abre mão ao fazer uma escolha.
\end{solution}
\vspace{4cm}

\question Considere a demanda $Q = 100 - 2P$.
\begin{parts}
  \part[1] Calcule a quantidade demandada quando $P = 10$.
  \begin{solution}$Q = 80$.\end{solution}
  \vspace{2.5cm}
  \part[2] Calcule a elasticidade-preço nesse ponto.
  \begin{solution}$\varepsilon = -2 \cdot 10/80 = -0{,}25$.\end{solution}
  \vspace{2.5cm}
\end{parts}

\question[1] Marque a alternativa correta: um bem de Giffen é aquele cuja demanda
\begin{checkboxes}
  \choice cai quando a renda sobe.
  \CorrectChoice sobe quando o próprio preço sobe.
  \choice não depende do preço.
  \choice sobe quando o preço do substituto cai.
\end{checkboxes}

\question[2] Verdadeiro ou falso? Justifique: ``Todo bem inferior é um bem de Giffen.''
\begin{solution}
  Falso. Todo bem de Giffen é inferior, mas não o contrário.
\end{solution}
\vspace{4cm}

\bonusquestion[1] Questão extra: explique o paradoxo da água e do diamante.
\vspace{3cm}

\end{questions}

\end{document}
